%% ATOM-IDS -- IETE Journal of Research submission manuscript
%% Layout follows the official IETE Journal of Research template supplied in
%% REf/IETE-Journal-of-Research-Template.pdf.
%% Compile: pdflatex ATOM-IDS-IETE -> bibtex ATOM-IDS-IETE -> pdflatex x2

\documentclass[10pt,letterpaper,twocolumn]{article}

\usepackage{cmap}
\usepackage[T1]{fontenc}
\usepackage[utf8]{inputenc}
\usepackage{newtxtext}
\usepackage[top=0.68in,bottom=0.70in,left=0.67in,right=0.67in,
            columnsep=0.24in,headheight=13pt]{geometry}
\usepackage{amsmath}
\usepackage{newtxmath}
\usepackage{graphicx}
\usepackage{booktabs}
\usepackage{multirow}
\usepackage{makecell}
\usepackage{siunitx}
\usepackage[ruled,vlined]{algorithm2e}
\usepackage{url}
\usepackage[hidelinks]{hyperref}
\usepackage{xurl}
\usepackage{cite}
\usepackage{titlesec}
\usepackage{fancyhdr}
\usepackage{caption}
\usepackage{dblfloatfix}

\setlength{\emergencystretch}{2em}
\setlength{\parindent}{1em}
\setlength{\parskip}{0pt}
\setlength{\textfloatsep}{7pt plus 2pt minus 2pt}
\setlength{\floatsep}{6pt plus 2pt minus 2pt}
\setlength{\intextsep}{6pt plus 2pt minus 2pt}
\captionsetup{font=small,labelsep=period,justification=justified,
              singlelinecheck=false}
\titleformat{\section}{\bfseries\normalsize}{\thesection.}{0.45em}{\MakeUppercase}
\titleformat{\subsection}{\bfseries\normalsize}{\thesubsection}{0.45em}{}
\titleformat{\subsubsection}{\itshape\normalsize}{\thesubsubsection}{0.45em}{}
\titlespacing*{\section}{0pt}{8pt plus 2pt minus 1pt}{3pt}
\titlespacing*{\subsection}{0pt}{6pt plus 2pt minus 1pt}{2pt}
\titlespacing*{\subsubsection}{0pt}{5pt plus 1pt minus 1pt}{2pt}
\pagestyle{fancy}
\fancyhf{}
\fancyhead[L]{\scriptsize ATOM-IDS}
\fancyhead[R]{\scriptsize\thepage}
\renewcommand{\headrulewidth}{0pt}
\renewcommand{\footrulewidth}{0pt}
\makeatletter
\renewcommand\@biblabel[1]{#1.}
\makeatother

\begin{document}

\twocolumn[{%
\begin{@twocolumnfalse}
\begin{center}
{\LARGE\bfseries ATOM-IDS: An Accumulator-Based Tiny On-Device Model for
Intrusion Detection in Microcontroller-Class IoT Devices\par}
\vspace{7pt}
{\large Ashish Kalra and Manoj Kumar\par}
\vspace{5pt}
{\small Department of Computer Science \& Engineering, Delhi Technological
University, Main Bawana Road, Delhi 110042, India\par}
{\small Corresponding author: Ashish Kalra
(\texttt{ashishkalra\_2k22phdco502@dtu.ac.in})\par}
\end{center}
\vspace{5pt}
\noindent\textbf{ABSTRACT}\par
\small
IoT intrusion detection requires high predictive accuracy while operating
within the storage, latency and energy limits of embedded devices. We present
ATOM-IDS, an Accumulator-based Tiny On-device Model for IDS, with a novel
fixed-shape grouped accumulator that replaces per-feature tokenisation. Each
feature is encoded by a binned embedding and a rank-1 component, assigned to
one of \(K\) fixed random groups, rescaled by a power-of-two factor and refined
by a shared residual MLP. The method is evaluated on Edge-IIoTset and
CIC-IIoT-2025 using a leakage-controlled, class-stratified chronological
protocol, and on CICIDS2017
as an external enterprise-traffic robustness benchmark. The model is deployed
on an ESP32-S3 edge device. With $d = 8$, $K = 8$, INT8 weights and FP32
arithmetic, ATOM-IDS achieves macro-F1 scores of 0.8007, 0.5638 and 0.6163 on
Edge-IIoTset, CIC-IIoT-2025 and CICIDS2017, respectively, with corresponding
ESP32-S3 inference latencies of \SI{243.5}{\micro\second},
\SI{250.6}{\micro\second} and \SI{277.4}{\micro\second}. The Edge-IIoTset build
consumes a measured differential energy of \SI{9.56}{\micro\joule} per
inference. It is $2.4\times$ faster and 7.4 percentage points higher in macro-F1
than the identically deployed MLP ($h = 64$) baseline on Edge-IIoTset.
Under the same protocol, the tested FT-Transformer requires substantially more
arithmetic and activation memory and does not fit the target deployment
budget. These results show that streaming grouped accumulation can provide
competitive intrusion detection within microcontroller resource limits.\par
\vspace{5pt}
\noindent\textbf{Keywords:} Edge computing, intrusion detection, IoT security,
on-device inference, quantisation, tabular deep learning, TinyML.
\vspace{10pt}
\end{@twocolumnfalse}
}]

\section{Introduction}
\label{sec:intro}

The rapid expansion of Internet of Things (IoT) deployments has enlarged the
network attack surface while many endpoints remain limited in memory,
computation and energy \cite{ferrag2020dl}. An intrusion detector running on
the endpoint can respond without transmitting every flow to a gateway,
reducing communication delay, retaining raw traffic locally and remaining
available during network disruption. In this study, kilobyte-scale model
storage and sub-millisecond inference are adopted as target constraints for a
single-core microcontroller.

Recent deep-learning IDS studies commonly use recurrent, attention-based or
hybrid models, often combined with metaheuristic feature selection and
hyperparameter optimisation. Most are evaluated on desktop or simulated platforms and do not
report flash consumption, working memory, measured latency or energy on
microcontrollers. TinyML studies demonstrate that intrusion detection can
operate on constrained hardware \cite{hamdouchi2025tinyml}, while hardware-aware
neural architecture search confirms the feasibility of device-constrained
models \cite{ortizgarces2026nas}. However, these approaches reuse
established computational structures instead of reconsidering how tabular
features are represented and aggregated.

A representative purpose-built lightweight alternative reduces transformer
cost through low-rank projection, separable convolution and quantisation
\cite{wang2024theseus}, while retaining per-feature tokenisation. For $N$ input
fields, tokenisation materialises an $N \times d$ intermediate and
self-attention operates on a sequence of length $L=N+1$ at $O(L^{2})$ cost
\cite{vaswani2017attention}. For tabular flows, feature
order is defined primarily by the data schema rather than a meaningful temporal
sequence. Consequently, token and attention costs may grow without a
deployment benefit. Moreover, analytic multiply--accumulate
counts alone cannot establish deployability because operator choice affects
wall-clock latency.

Evaluation rigor matters. Socket identifiers can act as label shortcuts,
and random splits can place related campaign records in both training and test
partitions, producing optimistic estimates \cite{arp2022dos}. A credible
on-device evaluation should therefore control identifier leakage, preserve
chronological order within each class and measure the exported model on the
target hardware.

To address these limitations, we propose ATOM-IDS, an Accumulator-based Tiny
On-device Model for IDS. Each feature is represented by a binned-embedding
lookup \cite{gorishniy2022embeddings} fused with a rank-1 component. The
resulting encodings are assigned to $K$ fixed random groups, accumulated in
streaming order, rescaled by a power-of-two factor and refined by a
weight-shared residual multilayer perceptron. This construction avoids
materialising an $N \times d$ token tensor; peak activation is instead a fixed
$K \times d$ group state independent of the number of input features.

ATOM-IDS is evaluated on Edge-IIoTset \cite{ferrag2022edgeiiotset},
CIC-IIoT-2025 \cite{firouzi2025datasense} and CICIDS2017, with the latter used
as an external enterprise-traffic robustness benchmark. Under the
leakage-controlled temporal protocol, the deployed $d=8$, $K=8$ model achieves
macro-F1 scores of 0.8007, 0.5638 and 0.6163 with ESP32-S3 latencies of
\SI{243.5}{\micro\second}, \SI{250.6}{\micro\second} and
\SI{277.4}{\micro\second}, respectively. On Edge-IIoTset, it consumes
\SI{9.56}{\micro\joule} per inference and is $2.4\times$ faster and 7.4
macro-F1 percentage points higher than the identically deployed MLP ($h=64$)
baseline.

The contributions of this study are:

\begin{itemize}
\item A fixed-shape grouped accumulator that combines binned and rank-1 feature
encodings with power-of-two rescaling and shared residual refinement,
eliminating per-feature token storage.

\item A leakage-controlled evaluation across three datasets that removes
socket identifiers, uses within-class chronological partitions and distinguishes
external robustness from cross-dataset transfer.

\item Controlled MLP, classical and transformer comparisons, together with
encoder and aggregation ablations, that locate the sources of accuracy and
compute differences.

\item Physical ESP32-S3 deployment reporting macro-F1, minority-class recall,
flash, working memory, measured latency, differential energy and device--host
parity over 1000 flows per firmware build.
\end{itemize}

Section~\ref{sec:related} reviews related work, Section~\ref{sec:method}
presents the methodology, Section~\ref{sec:results} discusses
experimental results and Section~\ref{sec:conclusion} concludes the paper.

\section{Literature review}
\label{sec:related}

Related work is organised around three questions: what recent IoT and wireless
sensor network (WSN) detectors model, whether lightweight designs are verified
on embedded hardware, and whether their evaluations support fair comparison.

\subsection{Learning-based IDS for IoT networks}
\label{sec:related:learning}

Six recent \emph{IETE Journal of Research} papers illustrate the prevailing
accuracy-oriented approach. Misbha et al. combine embedding, convolution and
bidirectional recurrence for DDoS and man-in-the-middle detection on KDDCup99
and UNSW-NB15 \cite{misbha2024security}. Grandhi et al. place a convolutional
GRU within an authenticated SDN--IoT framework and evaluate it on IoT-23
\cite{grandhi2025secure}. Yadav and Kumar join a stacked autoencoder, modified
ANFIS and metaheuristic optimisation on CICIDS2017 and UNSW-NB15
\cite{yadav2026anfis}. For WSNs, Ram and Saminathan combine preprocessing,
feature selection and an optimised self-attention progressive GAN on WSN-DS
\cite{ram2025sapgan}; Hemalatha et al. use filtering, feature selection and an
interpretable additive neural model on AWID \cite{hemalatha2025igann}; and
Divya and Kumar evaluate an optimised cross-stage partial recurrent network on
four established datasets \cite{divya2026crossstage}. These systems broaden
model choice and attack coverage, but use different datasets, splits and
metrics. None of the six reports the complete combination of microcontroller
flash, working memory, latency, energy and host--device output parity needed to
establish on-device feasibility.

Misbha et al. and Grandhi et al. emphasise learned spatial--temporal patterns;
Yadav and Kumar combine representation learning with fuzzy inference; and the
three WSN studies add optimisation around preprocessing, feature selection or
classifier training. Consequently, their headline values cannot be averaged
or ranked. They are used here to identify architectural trends and missing
deployment evidence rather than to claim direct predictive superiority.

\subsection{Lightweight and embedded IDS}
\label{sec:related:lightweight}

Model compression does not by itself demonstrate deployability. Wang et al.
use knowledge distillation to compress a transformer-derived detector to as
few as 788 parameters, but do not report microcontroller execution
\cite{wang2024theseus}. FlowTransformer separates input encoding, transformer
body and classification head; its experiments show that the head affects both
predictive performance and model size \cite{manocchio2024flowtransformer}.
Both studies motivate a transformer comparator, while leaving activation
memory and device measurements open.

TinyML studies supply stronger hardware evidence. Hamdouchi and Idri compare
six classifiers on an Arduino Uno and jointly report kappa, latency, static RAM
and flash \cite{hamdouchi2025tinyml}. Ortiz-Garc\'{e}s et al. apply
hardware-aware neural architecture search to feature-aligned IoT corpora and
report \SI{2.48}{\milli\second} latency, less than 20\,KB RAM and
\SI{0.39}{\milli\joule} per inference on Cortex-M4/ESP32 targets
\cite{ortizgarces2026nas}. Table~\ref{tab:related-compare} separates metric,
storage or RAM, latency and energy; NR denotes unreported device evidence, not
zero cost.

These embedded studies expose two distinct notions of ``lightweight.'' A small
parameter count controls persistent storage, whereas feasible execution also
depends on peak activation memory, operator support and arithmetic cost.
Hamdouchi and Idri provide unusually detailed resource measurements, but their
best reported point uses a different dataset and model family from the neural
comparators. Ortiz-Garc\'{e}s et al. demonstrate hardware-aware search and
energy measurement, although their feature-alignment and evaluation protocol
differ from those used here. ATOM-IDS therefore does not compare isolated
headline scores across papers. It reports external studies as deployment
context, then makes causal architecture comparisons only among models trained,
quantised and timed under the same local protocol and device conditions.

\subsection{Tabular representation learning and evaluation integrity}
\label{sec.related.tabular}

Network-flow records are tabular. FT-Transformer represents every feature as a
token and applies self-attention \cite{gorishniy2021revisiting}; numerical
feature embeddings can also improve non-transformer backbones
\cite{gorishniy2022embeddings}. Nevertheless, tree ensembles remain strong
baselines on medium-sized tabular datasets \cite{grinsztajn2022why}. ATOM-IDS
therefore retains binned feature encoding but replaces the feature-length token
sequence with fixed-shape grouped accumulation.

Evaluation design limits cross-paper conclusions. IDS corpora lack a standard
feature set, complicating comparison \cite{sarhan2022standard}; temporal and
spatial bias can inflate security-classification results
\cite{pendlebury2019tesseract}; and broader audits document recurring leakage,
sampling and metric errors in security machine learning \cite{arp2022dos}.
Accordingly, ATOM-IDS removes identifiers, fits preprocessing on training data,
uses temporal splits and reports macro-F1. CICIDS2017 is treated as a separate
enterprise-traffic robustness check, not evidence of unseen-dataset transfer.

\begin{table*}[!t]
\centering
\caption{Comparison of reviewed IDS evidence. Values are self-reported and are
not directly comparable across datasets or protocols. NR: not reported.}
\label{tab:related-compare}
\scriptsize\setlength{\tabcolsep}{2.2pt}
\renewcommand{\arraystretch}{1.10}
\begin{tabular}{@{}p{2.15cm}p{2.45cm}p{2.15cm}p{1.65cm}p{1.55cm}p{1.85cm}p{1.55cm}@{}}
\toprule
Study & Dataset(s) & Reported result & Hardware & Latency & Storage/RAM & Energy \\
\midrule
Misbha et al.\ \cite{misbha2024security} & KDDCup99, UNSW-NB15 & Classification metrics & NR & NR & NR & NR \\
Grandhi et al.\ \cite{grandhi2025secure} & IoT-23 & Accuracy & NR & NR & NR & NR \\
Yadav and Kumar \cite{yadav2026anfis} & CICIDS2017, UNSW-NB15 & Acc./P/R/F1 & NR & NR & NR & NR \\
Ram and Saminathan \cite{ram2025sapgan} & WSN-DS & Relative accuracy gain & NR & NR & NR & NR \\
Hemalatha et al.\ \cite{hemalatha2025igann} & AWID & DR/FAR improvement & NR & NR & NR & NR \\
Divya and Kumar \cite{divya2026crossstage} & Four IDS datasets & 99.91\% accuracy & NR & NR & NR & NR \\
Hamdouchi \& Idri \cite{hamdouchi2025tinyml} & NF-BoT-IoT-v2 & $\kappa=0.974$ & Arduino Uno & \SI{8}{\micro\second} & 4,730\,B SRAM & NR \\
Ortiz-Garc\'{e}s et al.\ \cite{ortizgarces2026nas} & TON\_IoT, IoT-23, Edge-IIoTset & macro-F1 $>0.95$ (IoT-23) & Cortex-M4/ESP32 & \SI{2.48}{\milli\second} & $<$20\,KB RAM & \SI{0.39}{\milli\joule} \\
\textbf{ATOM-IDS (this work)} & Edge-IIoTset & \textbf{macro-F1 $=0.8007$} & \textbf{ESP32-S3} & \textbf{\SI{243.5}{\micro\second}} & \textbf{6,603\,B read-only} & \textbf{\SI{9.56}{\micro\joule}} \\
\bottomrule
\end{tabular}
\end{table*}

The unresolved gap is an IDS representation whose working state does not grow
with the feature count and whose predictive and physical costs are measured
together. ATOM-IDS addresses that gap through grouped accumulation, matched
baselines, leakage-controlled temporal evaluation and ESP32-S3 parity,
latency, flash and energy measurements for this study.

\iffalse

This section examines and compares recent research addressing intrusion detection in IoT networks. The selected period provides a basis for examining emerging approaches, methodological improvements, and notable research directions in this domain. This section reviews existing research on intrusion detection in IoT networks, covering learning-based detectors, lightweight and transformer-based architectures, TinyML deployment on embedded hardware and tabular representation learning. The review highlights key performance outcomes and limitations while establishing a systematic framework for comparing the different approaches.

\subsection{Learning-based IDS for IoT networks}
\label{sec:related:learning}

In 2025, Villafranca and Cano \cite{villafranca2025edgedl} proposed a
lightweight, dataset-agnostic DNN architecture (256--128--64--Softmax) for IoT
intrusion detection in sustainable smart-agriculture, curating a benchmark of
18 IoT-IDS datasets and reporting 99.14\% average accuracy, including perfect
scores on N-BaIoT, Car-Hacking and CIC-IoT2022, with only around 1.2 million
parameters and no per-dataset tuning. In 2025, Okey et al.\
\cite{okey2025ipasecoit} proposed iPASecIoT, a single-model framework that
concurrently identifies IoT devices and detects intrusions from fine-grained
behavioural fingerprints, combining machine and deep learning with a modified
firefly algorithm under a kappa-score-based voting mechanism for adaptive
feature selection, reporting mean F1 scores above 98\% for both device
identification and threat classification across the CICIoMT2024, CICIoT2023
and UNSW2019 datasets.

In 2025, Catillo et al.\ \cite{catillo2025multicids} proposed MultiCIDS, which
treats IoT and cyber-physical-system monitoring data as multivariate time
series rather than isolated points, pipelining a semi-supervised autoencoder
that produces a per-point anomaly score with a sliding-window algorithm that
aggregates evidence into collective intrusion decisions, evaluated across
benchmark traces from IoT devices, a cyber-physical system and a ubiquitous
server. In 2026, Kilichev and Kim \cite{kilichev2026hybrid} proposed a
label-free hybrid framework that pairs an offline random-forest teacher with an
online Long Short-Term Memory student, combining probability calibration,
Mahalanobis-distance distribution monitoring and Page--Hinkley drift detection
under a leakage-resistant three-way protocol, reporting 0.9893 accuracy and a
mean latency of 12.97\,ms per sample on CICIoMT2024. These studies advance
accuracy, dataset diversity, or real-time adaptation to concept drift, but none
targets microcontroller-class execution or reports weight storage in
kilobytes and the reported inference timings, where given, are not measured
on a resource-constrained embedded device.

The common limitation in this group is therefore not a lack of model variety,
but the absence of a jointly reported accuracy--memory--latency trade-off
under a deployment-oriented split. ATOM-IDS is evaluated against that
trade-off rather than against accuracy alone.

Recent papers in the \emph{IETE Journal of Research} reflect the same emphasis
on learning-based network defence. Misbha et al.\ combine word embeddings,
convolution and bidirectional recurrence for IoT DDoS and man-in-the-middle
detection \cite{misbha2024security}; Grandhi et al.\ integrate a convolutional
GRU detector into an SDN--IoT security framework
\cite{grandhi2025secure}; and Yadav and Kumar modify an adaptive
neuro-fuzzy inference system for CICIDS2017 and UNSW-NB15
\cite{yadav2026anfis}. In wireless sensor networks, Ram and Saminathan use an
optimised self-attention progressive GAN \cite{ram2025sapgan}, Hemalatha et
al.\ propose an interpretable additive neural detector
\cite{hemalatha2025igann}, and Divya and Kumar combine feature selection with
a cross-stage partial recurrent network \cite{divya2026crossstage}. These six
studies motivate the section organisation adopted here and provide recent
journal-specific comparators. Their evaluations, however, do not report the
joint microcontroller measurements of flash, working memory, latency, energy,
and host--device parity reported for ATOM-IDS.

\subsection{Lightweight and transformer-based IDS}
\label{sec:related:lightweight}

In 2023, Wang et al.\ \cite{wang2023dynquant} applied dynamic quantisation to
compress an IoT intrusion detection model for edge execution while retaining
accuracy. The same line of work later used the BERT-of-Theseus module
replacement \cite{wang2024theseus}, substituting transformer blocks with
smaller successors during training so that compression and learning proceed
together. In 2024, Manocchio et al.\ \cite{manocchio2024flowtransformer}
released Flow Transformer, a modular framework in which the input encoding,
transformer body and classification head are interchangeable and showed that
the head dominates both accuracy and model size.

In 2025, Wu et al.\ \cite{wu2025tracer} proposed TRACER, an attack-aware
divide-and-conquer transformer for industrial IoT that routes traffic to
specialised sub-detectors. In 2025, Akuthota and Bhargava
\cite{akuthota2025transformer} reported a transformer-based IoT detector
competitive with recurrent baselines. These designs substantially reduce
parameter count, yet retain per-feature tokenisation and self-attention and
report latency on desktop hardware, so the compression demonstrated does not
establish microcontroller feasibility.

This distinction motivates a direct tokenising baseline in the experiments:
the comparison measures both predictive performance and the intermediate
activation state that must be resident on the device.

\subsection{TinyML and on-device intrusion detection}
\label{sec:related:tinyml}

In 2025, Hamdouchi and Idri \cite{hamdouchi2025tinyml} compared three singular
classifiers (decision tree, Gaussian na\"{\i}ve Bayes, multilayer perceptron)
against three tree ensembles (random forest, XGBoost, extra trees) on an
Arduino Uno using the NF-ToN-IoT-v2 and NF-BoT-IoT-v2 datasets, ranking them by
the Scott-Knott test and Borda count voting, with Cohen's kappa, latency,
static RAM and flash memory. Their overall conclusion favoured singular
models, with a 20-feature, 56-unit multilayer perceptron the best model on
NF-ToN-IoT-v2 ($\kappa = 0.935$, 7,010 bytes of static RAM,
\SI{9.6}{\micro\second} latency). the single fastest and smallest result
across both datasets, however, was a 13-feature extra-trees ensemble on
NF-BoT-IoT-v2, at $\kappa = 0.974$, 4,730 bytes of static RAM and
\SI{8}{\micro\second} latency. In 2025, Amuthadevi et al.\
\cite{amuthadevi2025graboost} proposed GraBoost-AAGT, coupling adaptive
artificial gorilla troops optimisation with extreme gradient boosting after
Z-score normalisation and Fast Fourier Transform-based feature extraction,
reporting 99.50\% accuracy with reduced CPU and energy overhead.

In 2025, Fusco et al.\ \cite{fusco2025siamese} combined Siamese neural networks
with gradient sparsification for few-shot federated training on microcontroller
units, selecting the sparsification ratio each round by an epsilon-greedy
reinforcement learning policy. In 2026, Ortiz-Garcés et al.\
\cite{ortizgarces2026nas} applied hardware-aware Neural Architecture Search
over a feature-aligned corpus drawn from TON\_IoT, IoT-23 and Edge-IIoTset,
compressing it to 32--64 features, retaining over 99\% of the informative
variance and validated the model at \SI{2.48}{\milli\second} latency, under
20\,KB of RAM and \SI{0.39}{\milli\joule} per inference. The wider TinyML
systems literature explains why such measurement is necessary. TensorFlow Lite
Micro \cite{david2021tflm} shows that interpreter and kernel design dominate
execution time, MCUNet \cite{lin2020mcunet} shows that architecture and memory
scheduling must be co-designed under a fixed SRAM budget and MLPerf Tiny
\cite{banbury2021mlperftiny} standardises embedded latency and energy
measurement. These systems confirm that on-device detection is achievable, but
each deploys an established model family rather than reconsidering the operator
itself.

ATOM-IDS follows the same measurement discipline while changing the operator
that builds the representation: feature encodings are accumulated into a fixed
number of groups, so the activation state is independent of the feature count.

\subsection{Tabular representation learning and evaluation integrity}
\label{sec.related.tabular}

Because a flow record is tabular rather than sequential, representation has
been studied independently of intrusion detection. FT-Transformer
\cite{gorishniy2021revisiting} embeds each feature as a token and applies
self-attention across the sequence, while TabNet \cite{arik2021tabnet} applies
sequential attention masks at each decision step. both materialise a
per-feature intermediate representation that grows with the feature count.
Grinsztajn et al.\ \cite{grinsztajn2022why} found tree-based models still
outperform deep networks on typical tabular data and Shwartz-Ziv and Armon
\cite{shwartzziv2022tabular} reached the same conclusion across a range of
benchmarks, which is why a gradient-boosted arm is retained as a reference
here. Gorishniy et al.\ \cite{gorishniy2022embeddings} showed that
quantile-binned embeddings of numerical features substantially improve tabular
networks, a result that the proposed encoder adopts.

Benchmarking practice is a further concern. Sarhan et al.\
\cite{sarhan2022standard} observed that IDS datasets lack a common feature set
and that socket identifiers can leak the label almost directly. Pendlebury et
al.\ \cite{pendlebury2019tesseract} showed that random splits let a classifier
memorise campaign artefacts that will not recur in deployment. Arp et al.\
\cite{arp2022dos} catalogued these and related pitfalls across security machine
learning. Shafi and Habibi Lashkari \cite{shafi2025iotzwave} address a related
gap from the dataset side, profiling smart-home devices and characterising
Z-Wave traffic and intruder behaviour to build a large-scale, multi-protocol
IoT intrusion detection dataset, which underlines how much a reported
accuracy depends on the dataset it was measured against. Accuracy reported
without leakage control and a temporal split, therefore, does not characterise
field behaviour.

The present study consequently reports temporal within-dataset results for all
three corpora and labels CICIDS2017 explicitly as an external enterprise
traffic check. It does not claim unseen-campaign transfer because the model is
trained and tested separately within each corpus.


Table~\ref{tab:related-compare} positions ATOM-IDS against the closest
deployed or dataset-scale comparators from the reviewed literature, each on its
own self-reported metric. 

\begin{table}[!ht]
\centering
\caption{ATOM-IDS against the closest deployed or dataset-scale comparators
from the reviewed literature.}
\label{tab:related-compare}
\footnotesize\setlength{\tabcolsep}{3pt}
\renewcommand{\arraystretch}{1.15}
\begin{tabular}{@{}p{2.0cm}p{2.4cm}p{1.6cm}p{2.0cm}p{1.55cm}p{1.95cm}@{}}
\toprule
Study & Dataset(s) & Hardware & Deployed metric & Latency & Memory \\
\midrule
Hamdouchi \& Idri \cite{hamdouchi2025tinyml} &
  NF-ToN/BoT-IoT-v2 & Arduino Uno &
  $\kappa = 0.974$ & \SI{8}{\micro\second} & 4,730\,B SRAM \\
Ortiz-Garc\'{e}s et al.\ \cite{ortizgarces2026nas} &
  TON\_IoT, IoT-23, Edge-IIoTset & Cortex-M4 / ESP32 &
  --- & \SI{2.48}{\milli\second} & $<$20\,KB RAM, 0.39\,mJ/inf. \\
Villafranca \& Cano \cite{villafranca2025edgedl} &
  18 IoT-IDS datasets & not deployed &
  99.14\% avg.\ acc.\ & --- & $\sim$1.2M params \\
\textbf{ATOM-IDS (this work)} &
  Edge-IIoTset & \textbf{ESP32-S3} &
  \textbf{macro-F1 $=$ 0.8007} & \textbf{\SI{243.5}{\micro\second}} &
  \textbf{4,912\,B flash, 0.0096 mJ/inf.} \\
\bottomrule
\end{tabular}
\end{table}

To address these limitations, this work proposes ATOM-IDS, an accumulator-based,
hardware-aware tabular model that removes tokenisation in favour of a fixed-shape grouped
accumulator, sizes its activation state by the group count rather than the
feature count, is evaluated under a leakage-controlled temporal protocol, and
is verified for host-device output parity on an ESP32-S3 microcontroller.


\fi

\section{Proposed methodology}
\label{sec:method}

ATOM-IDS is a hardware-aware tabular classifier whose representation is built
by streaming feature encodings into a fixed number of groups. The forward path
comprises leakage-controlled preprocessing, a fused binned and rank-1 encoder,
balanced fixed grouping, power-of-two rescaling, a shared residual perceptron
and a linear classifier. Figure~\ref{fig:arch} shows this path and
Algorithm~\ref{alg:atom} states the training and export procedure.

\begin{figure}[!t]
\centering
\includegraphics[width=0.96\linewidth,height=0.67\textheight,keepaspectratio]{figs/fig1_architecture.pdf}
\caption{Architecture of the proposed ATOM-IDS framework.}
\label{fig:arch}
\end{figure}

\subsection{Dataset description and preprocessing}
\label{sec:method:data}

Three datasets test the architecture under different schemas.
Edge-IIoTset \cite{ferrag2022edgeiiotset} contains 2.22 million flow records
from an IoT/IIoT testbed, with benign traffic and 14 attack classes.
CIC-IIoT-2025 \cite{firouzi2025datasense} provides per-device window
aggregates; this study uses its one-second windows and eight-class categorical
labels. CICIDS2017 \cite{sharafaldin2018cicids} contains 2.83 million
bidirectional enterprise flows collected over five days and 15 classes. It is
used as an external enterprise-traffic robustness benchmark, not as an IoT
dataset or an unseen-dataset transfer test. Every model is trained and tested
separately within each corpus. Table~\ref{tab:datasets} gives the retained
feature count $N$, class count $C$ and partition sizes.

\begin{table}[!ht]
\centering
\caption{Dataset description after leakage control.}
\label{tab:datasets}
\begin{tabular}{lrrrrr}
\toprule
Dataset & $N$ & $C$ & Train & Val & Test \\
\midrule
Edge-IIoTset  & 54 & 15 & 385,478 & 120,669 & 117,959 \\
CIC-IIoT-2025 & 71 &  8 &  78,851 &  26,511 &  23,230 \\
CICIDS2017    & 70 & 15 & 201,385 &  66,851 &  67,404 \\
\bottomrule
\end{tabular}
\end{table}

Preprocessing follows one discipline across the three corpora. First, direct
identifiers and label shortcuts are removed. These include flow identifiers,
addresses, ports, payload-like strings, URIs, DNS names, MQTT topics and any
binary attack flag that directly reveals the multiclass target. CICIDS2017
retains its numeric protocol field as a categorical feature. CIC-IIoT-2025
retains numeric count summaries but not the raw identifier-list fields. This
restriction follows security-ML guidance on shortcut leakage
\cite{arp2022dos,sarhan2022standard}.

Second, a class-stratified chronological split assigns the earliest 60\% of
each class to training, the next 20\% to validation and the final 20\% to
testing. Edge-IIoTset uses frame time, CIC-IIoT-2025 uses window start and
CICIDS2017 uses flow time with capture order breaking timestamp ties. For
CICIDS2017, rank within each class is used because minute-resolution ties can
otherwise collapse a short attack burst into one partition. This construction
keeps every class represented and preserves order \emph{within} each class. It
is not a globally disjoint time-window evaluation: timestamps from different
classes may overlap across partitions. It therefore measures within-campaign,
class-stratified temporal generalisation and makes no claim about an unseen
campaign or attack family \cite{pendlebury2019tesseract}.

Third, exact repetitions are controlled before sampling. The duplicate key is
the retained raw feature set together with the class label; the first record in
source order is kept and later occurrences of that same feature-label tuple
are removed. Thus differently labelled collisions are retained. This operation
removes 14.0\% of Edge-IIoTset rows and 35.1\% of CIC-IIoT-2025 windows. It
does not fit a statistic, alter a value or move a record between temporal
partitions. Within each partition and class, records above the stated budget
are sampled without replacement using preprocessing seed 0. The respective
train/validation/test caps are 60,000/20,000/20,000 for Edge-IIoTset,
30,000/10,000/10,000 for CIC-IIoT-2025 and 60,000/20,000/20,000 for
CICIDS2017.

All fitted transforms use training data only. Numeric parse failures and
non-finite training medians are replaced by the training median or zero when
that median is itself non-finite. Edge-IIoTset categorical vocabularies contain
the 15 most frequent training levels plus an \texttt{OTHER} indicator, which
also receives unseen validation and test levels. Training-constant columns are
dropped. Non-negative heavy-tailed lengths, sizes, checksums, sequence fields
and counts are clipped below at zero and transformed by $\log(1+x)$. Remaining
features are z-scored with training means and standard deviations and clipped
to $[-8,8]$. The resulting vocabularies, medians, scaling statistics and
feature order are reused unchanged for validation, testing and firmware export.

\subsection{Fused binned-embedding and rank-1 encoder}
\label{sec:method:encoder}

Let $x\in\mathbb{R}^{N}$ be a standardised record. For feature $i$, seven
training-set empirical quantiles define $B=8$ bins; repeated quantiles are
retained, so tied or near-constant features may have fewer than eight effective
intervals. If $b_i(x_i)$ is the number of edges strictly below $x_i$, the
feature encoding is

\begin{equation}
\label{eq:encoder}
e_i=E[i,b_i(x_i)]+x_iV[i],
\end{equation}

where $E\in\mathbb{R}^{N\times B\times d}$ is a learned lookup table and
$V\in\mathbb{R}^{N\times d}$ supplies a rank-1 linear direction. The lookup
provides piecewise-constant nonlinearity without multiplications, while the
rank-1 term preserves within-bin variation at $d$ multiply-accumulates (MACs)
per feature. This design adopts the demonstrated value of numerical feature
embeddings \cite{gorishniy2022embeddings} without retaining a feature-length
token sequence after encoding.

\subsection{Balanced fixed grouping}
\label{sec:method:grouping}

The feature indices are assigned to $K$ groups by shuffling the repeating
sequence $0,1,\ldots,K-1$ with topology seed 0. Group sizes therefore differ by
at most one and no group is empty when $N\geq K$. This single balanced map is
fixed for every dimension, training seed and quantisation run within a corpus;
it is not selected by validation performance. The group state is

\begin{equation}
\label{eq:group}
G_k=\sum_{i:g(i)=k}e_i.
\end{equation}

The implementation orders features so that group members are contiguous and
streams each encoding directly into its destination accumulator. It never
materialises the $N\times d$ encoding stack. The live group state is therefore
$K\times d$ elements, independent of $N$, although the input vector, bin edges
and stored encoder remain linear in $N$. The deployed kernel holds the
accumulator in float32, so its byte cost is four times its element count.

The deployed model is the $R=1$ case of an earlier adaptive design. Its
$R=4$ control contains four candidate assignment matrices and a width-8 router
that selects a codebook entry for each input. That control is retained only for
the matched ablation in Section~\ref{sec:ablation}; deployment uses the single
fixed map above. Because the production map is constant, variation across the
reported training seeds measures optimisation stability, not sensitivity to
the topology. Independently redrawn balanced maps were not evaluated and are
listed as future work.

\subsection{Power-of-two rescaling}
\label{sec:method:pow2}

Training learns one unrestricted scale $\alpha_k$ per group. At export, it is
converted to a signed integer exponent

\begin{equation}
\label{eq:pow2}
q_k=\mathrm{round}\!\left(\log_2(|\alpha_k|+10^{-8})\right),\qquad
G_k\leftarrow 2^{q_k}G_k.
\end{equation}

Thus the firmware stores $K$ exponents rather than $K$ floating-point scales.
The export-time rounding is deterministic and no gradient through
$\mathrm{round}(\cdot)$ is claimed. In a future integer-only kernel the factor
$2^{q_k}$ can be implemented as a shift; in the measured float32 kernel it is
an exactly representable binary scaling. Predictive effects of this constraint
belong to the quantisation ablation in Section~\ref{sec:quant}, rather than to
the method definition.

\subsection{Residual refinement and classification}
\label{sec:method:head}

A two-layer residual perceptron is shared across all groups:

\begin{equation}
\label{eq:refine}
G_k\leftarrow G_k+\gamma\left[W_2\,\mathrm{ReLU}(W_1G_k+b_1)+b_2\right],
\end{equation}

where $W_1\in\mathbb{R}^{h\times d}$,
$W_2\in\mathbb{R}^{d\times h}$, $h=\max(2,\lfloor d/2\rfloor)$, and
$\gamma$ is learned. The biases $b_1$ and $b_2$ are included in both training
and firmware. Sharing makes this refinement independent of $K$ in parameter
count. The refined groups are concatenated and mapped to class logits,

\begin{equation}
\label{eq:head}
z=W_c[G_1;\ldots;G_K]+b_c,\qquad \hat y=\arg\max_c z_c.
\end{equation}

Concatenation is necessary because direct mean pooling would make every fixed
partition identical to the mean of all feature encodings and would discard
which group received each feature.

\subsection{Computational and storage cost}
\label{sec:method:cost}

Ignoring bias additions and nonlinearities, the analytic MAC count is

\begin{equation}
\label{eq:macs}
\mathrm{MACs}=Nd+K(2dh)+KdC.
\end{equation}

The encoder lookup contributes no MACs. Group formation additionally requires
$Nd$ additions, binning requires at most $N(B-1)$ ordered comparisons, and
rescaling requires $Kd$ binary scalings in the measured kernel. These
operations are included in device latency even though they are not counted as
MACs. There is no attention term quadratic in $N$.

Storage and working memory are reported separately. The INT8 weight count
contains $E$, $V$, $W_1$, $W_2$ and $W_c$. The complete read-only artifact also
contains float biases and scales, the gate, quantile edges, permutation, group
boundaries and shift exponents. Consequently, the ``INT8 weight bytes'' columns
in the deployment tables must not be interpreted as total firmware flash;
Table~\ref{tab:device-summary} reports the complete constant
footprint. Working memory includes the input, group accumulator, refined group
copy, hidden vector and output logits.

The accounting can be checked directly from the architecture. For
Edge-IIoTset at $d=8$, $K=8$ and $C=15$, the stored INT8 matrices contain
$N B d=3{,}456$ embedding entries, $N d=432$ rank-1 entries, 64 residual-layer
weights and $K d C=960$ classifier weights, totalling 4,912 weight bytes. The
trained model additionally contains 27 biases, the residual gate and eight
group scales, giving 4,948 learned scalar parameters. At export, the group
scales become eight shift exponents, while the full header adds the bin edges,
permutation, group boundaries, float scales and biases. This reconciliation is
used for every configuration and prevents parameter count, quantised weight
storage and total read-only storage from being treated as interchangeable.

Feature count affects storage more strongly than arithmetic in this design.
The $E$ and $V$ tensors grow linearly with $N$, whereas the group state remains
$K d$ and the classifier depends on $K d C$. Consequently, feature selection
would primarily reduce the embedding table and bin-edge footprint. It removes
only the corresponding $d$ encoder MACs per feature and does not alter the
residual or classifier terms. This distinction matters when comparing ATOM-IDS
with TinyML studies that use feature selection as their main compute-reduction
mechanism \cite{hamdouchi2025tinyml}.

\begin{algorithm}[!t]
\caption{ATOM training and on-device evaluation.}
\label{alg:atom}
\KwIn{Raw records, topology seed 0, training seed $s$, $d$, $K$, $B=8$.}
\KwOut{Macro-F1, per-class metrics, latency, memory, energy and parity.}
Remove identifiers and assign class-stratified chronological partitions\;
Within each class, retain the first occurrence of each feature-label tuple;
sample capped strata with seed 0\;
Fit imputation, encoding, scaling and bin edges on training data only\;
Construct one balanced group map with topology seed 0 and freeze it\;
Initialise $E,V,W_1,b_1,W_2,b_2,W_c,b_c,\alpha$ and $\gamma$ with seed $s$\;
\For{$\mathrm{epoch}=1$ \KwTo $25$}{
  Encode by \eqref{eq:encoder}, accumulate by \eqref{eq:group}, refine by
  \eqref{eq:refine} and classify by \eqref{eq:head}\;
  Step Adam on class-weighted cross-entropy and retain the best validation checkpoint\;
}
Apply PTQ or QAT; round group scales by \eqref{eq:pow2}; emit the C header\;
Flash the ESP32-S3 and verify 1,000 host-device class indices\;
Measure warm median latency, complete memory footprint and differential energy\;
\end{algorithm}

\subsection{Training, quantisation and device measurement}
\label{sec:method:training}

All neural models use Adam with learning rate $3\times10^{-3}$, batch size
1024 and 25 epochs. Inverse-frequency class weights are applied to
cross-entropy, and the checkpoint with the highest validation macro-F1 is
retained. Weighting changes the objective without synthesising records across
a temporal boundary, unlike applying an oversampler before splitting. The
FP32 frontier uses ten training seeds. These seeds control PyTorch and NumPy
initialisation and minibatch order; they do not redraw the fixed group map.

All neural comparators consume the same cached partitions and transformed
feature vectors. The MLP arms use the same optimiser, loss, epochs, batch size,
checkpoint rule and training seeds. The matched transformer additionally uses
the same binned-plus-rank-1 encoder, so the remaining architectural difference
is self-attention over $N+1$ tokens versus accumulation into $K$ groups.
Gradient-boosted trees are retained as a predictive reference but are neither
placed on the neural MAC axis nor deployed, because node traversal is not
equivalent to a dense MAC count. These rules keep causal comparisons within
the present experiment rather than across incompatible studies.

The quantisation study distinguishes PTQ, QAT and deployment. Weights are
stored with symmetric per-tensor INT8 scales. During QAT, unsigned asymmetric
8-bit fake quantisers with explicit zero points are applied at the input and
intermediate activation sites, while straight-through estimators pass gradients
through weight and activation rounding. Biases and the residual gate remain
float32. At export, the learned group scales are converted to the integer
exponents of \eqref{eq:pow2}. A C header then records INT8 weights, float
scales and biases, bin edges, the balanced permutation, group boundaries and
shift exponents.

The deployed implementation is accurately described as \emph{INT8-weight,
FP32-arithmetic}. It reads INT8 weights from flash, dequantises them with their
per-tensor scales and performs products, accumulation, activations and group
state storage in float32. Fake-quantised activation values used during QAT are
reproduced, but activations are not stored or accumulated as integers on the
device. Therefore, the reported weight-storage reduction is real, while the
measurements do not establish integer-only latency. An int8-by-int8/int32
kernel of the TensorFlow Lite Micro type \cite{david2021tflm} remains future
work.

Firmware is compiled with the Arduino ESP32 core for an ESP32-S3 at 240 MHz.
For every deployed build, 1,000 test records selected without replacement with
seed 0 are replayed and device class indices are checked against host indices.
All 15 builds attain 1,000/1,000 agreement. Latency is the median of 256 warm
inferences measured in CPU cycles; the first cache-cold call is reported
separately. Energy is obtained from the difference between active and idle
power multiplied by measured inference time, using the protocol of
Section~\ref{sec:energy}. Feature extraction from raw packets is outside the
timed region for every model.

Parity is defined on the predicted class index rather than exact logits. This
detects export errors that alter a decision while allowing harmless numerical
differences below the decision boundary. The 1,000 replayed records are a
deterministic verification sample, not another accuracy test; macro-F1 is
computed on the complete held-out test partition. Likewise, latency covers
classification of a precomputed feature vector and excludes packet capture,
flow assembly and feature extraction. The energy value is therefore
model-inference energy under the stated conditions, not the total energy of an
end-to-end network monitor.

\iffalse

This work proposes ATOM-IDS, an accumulator-based, hardware-aware tabular model
built around a novel fixed-shape grouped accumulator and designed to execute
within the storage and latency budget of a microcontroller-class endpoint. Raw
flow records are first pre-processed under a leakage-controlled
protocol, then encoded by a fused binned-embedding and rank-1 operator,
accumulated into a fixed number of random groups, rescaled by a power-of-two
shift, refined by a weight-shared residual multilayer perceptron, and
classified from the concatenated group state. Figure~\ref{fig:arch} shows the
forward pass, and Algorithm~\ref{alg:atom} summarises construction, training
and on-device evaluation.

\begin{figure}[!t]
\centering
\includegraphics[width=0.96\linewidth,height=0.67\textheight,keepaspectratio]{figs/fig1_architecture.pdf}
\caption{Architecture of the proposed ATOM-IDS framework.}
\label{fig:arch}
\end{figure}

\subsection{Dataset description and preprocessing}
\label{sec:method:data}

ATOM-IDS is evaluated on three intrusion-detection datasets with distinct data
schemas. Edge-IIoTset \cite{ferrag2022edgeiiotset} is a flow-based dataset
containing 2.22 million records across 15 classes, including 14 attack families
and benign traffic, collected from a purpose-built IoT/IIoT testbed.
CIC-IIoT-2025 \cite{firouzi2025datasense} provides per-device, per-time-window
feature aggregates, from which the 1-second window and 8-class categorical
labels are used in this study. CICIDS2017 \cite{sharafaldin2018cicids} is the
long-standing enterprise-capture benchmark for this field, comprising 2.83
million bidirectional flows over a five-day capture with 15 classes. It is
included as an external enterprise-traffic robustness benchmark, since it is
neither an IoT testbed nor a per-device aggregate and its wide prior use makes
the architecture's behaviour directly comparable with published work. The
experiments remain within-dataset temporal evaluations. CICIDS2017 is not
presented as an unseen-dataset transfer test.
Table~\ref{tab:datasets} summarises the three datasets after preprocessing,
where $N$ denotes the number of retained features and $C$ the number of
classes.

\begin{table}[!ht]
\centering
\caption{Dataset description after leakage control.}
\label{tab:datasets}
\begin{tabular}{lrrrrr}
\toprule
Dataset & $N$ & $C$ & Train & Val & Test \\
\midrule
Edge-IIoTset  & 54 & 15 & 385,478 & 120,669 & 117,959 \\
CIC-IIoT-2025 & 71 &  8 &  78,851 &  26,511 &  23,230 \\
CICIDS2017    & 70 & 15 & 201,385 &  66,851 &  67,404 \\
\bottomrule
\end{tabular}
\end{table}

All three datasets are processed under a single discipline intended to prevent
the label from leaking through identifying fields. First, socket identifiers,
comprising source and destination addresses and all port fields, are dropped
together with high-cardinality application-layer strings such as payloads,
URIs, DNS query names and MQTT topics, and with any binary attack flag that
acts as a direct shortcut to the multi-class label. On CICIDS2017 this removes
the flow identifier, both addresses and both port fields, while the numeric
protocol field is retained as a categorical feature. These fields identify
attacks almost directly \cite{arp2022dos} and account for much of the
implausibly high accuracy reported on these corpora \cite{sarhan2022standard}.
On CIC-IIoT-2025 the raw identifier list columns are removed while their
numeric count companions are retained. Second, the retained feature columns are
deduplicated globally, so that an identical feature vector cannot appear in two
partitions. this removes 14.0\% of Edge-IIoTset rows and 35.1\% of
CIC-IIoT-2025 windows. Because this filter runs before the partitions are
populated, its scope is stated precisely. It removes rows and never edits them,
the partition boundaries are fixed by timestamp independently of it, and every
fitted statistic in the third step below is computed on the training partition
after splitting, so no value, label or statistic from a validation or test
record enters training. Second, categorical fields on Edge-IIoTset are
one-hot encoded using the fifteen most frequent levels plus an indicator for
all others, heavy-tailed numeric fields such as lengths, sizes, checksums and
counts are compressed by a $\log(1+x)$ transform, and all numeric features are
standardised by z-scoring on training statistics and clipped to the range plus
or minus eight.

Each class is then split 60/20/20 by timestamp, using frame time on
Edge-IIoTset, window start on CIC-IIoT-2025 and flow timestamp on CICIDS2017,
so that the model never trains on the future of a campaign against which it is
later tested while every class remains present in every partition. On
CICIDS2017 the split is taken over the \emph{rank} of each record in
chronological order rather than over the timestamp value, with ties broken by
capture order. The distinction matters on that corpus and is recorded here
because it is easy to get wrong: its timestamps are minute-resolution and
several attacks are captured in a single short burst, so a class can have
identical 60th and 80th percentile timestamps, and a value-based comparison
then places every record of that class in one partition. Splitting on rank
keeps the partition chronological while guaranteeing the intended proportions
for every class regardless of timestamp granularity. a class left with no test
support would otherwise score zero F1 for every model and depress macro-F1 by
$1/C$ uniformly, which is a measurement artefact rather than a result. Because
all three datasets are dominated by a few classes, each split is additionally
capped per class, at 60,000/20,000/20,000 records on Edge-IIoTset,
30,000/10,000/10,000 on CIC-IIoT-2025 and 60,000/20,000/20,000 on CICIDS2017. Categorical
vocabularies, z-scoring statistics and quantile bin edges are fitted on the
training partition alone and reused unchanged for validation, test and device
export. This is a within-campaign temporal split. It measures temporal
generalisation and does not hold out an unseen campaign or attack family, a
boundary stated explicitly rather than left to be inferred
\cite{pendlebury2019tesseract}.

\subsection{Fused binned-embedding and rank-1 encoder}
\label{sec:method:encoder}

Let $x$ be the standardised feature vector of one flow, with $N$ components.
For each feature $i$, let $\mathrm{bin}(x_i)$ denote the index of the quantile
bin into which $x_i$ falls, where the bin edges are the $B-1$ interior
empirical quantiles of feature $i$ computed on the training partition, with
$B = 8$. The encoding of feature $i$ is

\begin{equation}
\label{eq:encoder}
e_{i} \;=\; E\!\left[i,\, \mathrm{bin}(x_{i})\right] \;+\; x_{i}\, V\!\left[i\right],
\end{equation}

where $E$ is a learned embedding table of shape $N \times B \times d$ and $V$
is a learned rank-1 direction of shape $N \times d$. The lookup term costs
storage but no multiplications, so only the rank-1 term contributes $d$
multiply-accumulate operations per feature. The binned term supplies the
nonlinearity that a pure rank-1 map lacks: with the rank-1 term alone the
entire feature path reduces to a structured-sparse linear map and accuracy
collapses towards logistic regression, which Section~\ref{sec:ablation}
quantifies. This encoder follows the established
evidence that quantile-binned embeddings materially improve tabular networks
\cite{gorishniy2022embeddings}.

\subsection{Grouped accumulation with fixed random topology}
\label{sec:method:grouping}

A frozen assignment $g$, mapping each of the $N$ feature indices to one of $K$
groups, is drawn once uniformly at random and never trained. The group state is
the sum of the encodings assigned to that group,

\begin{equation}
\label{eq:group}
G_{k} \;=\; \sum_{i \,:\, g(i)\,=\,k} e_{i}.
\end{equation}

The sum is formed by accumulating each encoding directly into its destination
group, so the $N \times d$ stack of encodings is never held in memory. Peak
activation state is therefore the $K \times d$ accumulator, whose element count
is independent of the feature count. only the input vector and the bin edges
remain of order $N$. This is a statement about the number of live elements
rather than about their runtime type: the deployed kernel of
Section~\ref{sec:method:training} holds those $K \times d$ elements as float32,
so the byte figures in Tables~\ref{tab:mem-edge} and \ref{tab:mem-cic} are the
element count multiplied by four, and an integer-only kernel would reduce the
constant without changing the order. On device the features are gathered into a
permuted order in which the members of a group are contiguous, so each group
accumulator stays resident across its whole range and no scattered writes are
performed. This structure is closer to a tied-weight, hard-routed mixture
\cite{shazeer2017moe} than to a transformer, but the routing is fixed rather
than input-conditional.

The fixed assignment used above is the $R = 1$ case of a more general operator,
and $R$ is defined here because the ablation of Section~\ref{sec:ablation}
varies it. Let the codebook $C_b$ be a tensor of shape $R \times N \times K$
holding $R$ candidate assignment matrices, each of whose $N$ rows is a one-hot
distribution over the $K$ groups, and let a small input-conditional router map
the standardised input through a two-layer network of width 8 to a one-hot
selection $r(x)$ over the $R$ entries. The effective assignment is then
$A(x) = \sum_{r} r_{r}(x)\, C_{b,r}$, and the group state of
\eqref{eq:group} becomes $G_{k} = \sum_{i} A(x)_{ik}\, e_{i}$. The parameter
$R$ is therefore the size of the topology codebook: the number of distinct
feature-to-group assignments the router may select between for a given input.
Setting $R = 1$ collapses $r(x)$ to a constant, removes the router, freezes $A$
to the single random map $g$, and recovers exactly the deployed model described
above. Setting $R > 1$ adds the learned codebook and input-conditional routing
at a cost of $8N + 8R$ further multiply-accumulate operations per inference,
plus the codebook parameters. ATOM-IDS deploys $R = 1$. the $R = 4$ variant is
trained only as the control against which that choice is tested in
Section~\ref{sec:ablation}.

\subsection{Power-of-two group requantisation}
\label{sec:method:pow2}

Each group is rescaled by a learned per-group factor constrained to the nearest
power of two,

\begin{equation}
\label{eq:pow2}
G_{k} \;\leftarrow\; 2^{\,\mathrm{round}(\log_{2} \alpha_{k})}\, G_{k},
\end{equation}

so that the rescale is exactly representable in binary and reduces, under a
pure-integer kernel, to an arithmetic shift rather than a general
multiplication. The constraint also removes $K$ freely learned per-group scales
from the parameter set. At $d = 8$ its accuracy cost relative to a free
per-group scale is within seed noise, at $+0.002$ macro-F1. at $d = 24$ it
costs 0.013 macro-F1 on Edge-IIoTset and 0.024 on CIC-IIoT-2025, consistently
across seeds. The constrained form is the one retained in the deployed INT8
path.

\subsection{Residual refinement and classification head}
\label{sec:method:head}

A single two-layer residual multilayer perceptron, with weights shared across
all $K$ groups, refines each group vector,

\begin{equation}
\label{eq:refine}
G_{k} \;\leftarrow\; G_{k} \;+\; \gamma\, W_{2}\, \mathrm{ReLU}(W_{1}\, G_{k}),
\end{equation}

where $W_{1}$ has shape $(d/2) \times d$, $W_{2}$ has shape $d \times (d/2)$,
and $\gamma$ is a scalar gate. Sharing the weights across groups keeps the
parameter cost independent of $K$. The $K$ refined vectors are then
concatenated and classified by a single linear layer,

\begin{equation}
\label{eq:head}
\hat{y} \;=\; \mathrm{softmax}\!\left(W_{c}\,[\,G_{1};\, \ldots;\, G_{K}\,] + b_{c}\right).
\end{equation}

The readout is concatenation rather than pooling, and this is a requirement
rather than a preference. For a readout applied directly to the grouped sums of
\eqref{eq:group}, the mean over groups equals the mean over all encodings for
every assignment $g$, so all groupings present the classifier with an identical
vector and the topology becomes unlearnable. Placing the pool after the rescale
and the refinement weakens this identity without repairing it, because the
outer sum still discards which group each feature was assigned to.

\subsection{Computational cost}
\label{sec:method:cost}

The analytic multiply-accumulate count per inference is the sum of the rank-1
encoder, the shared residual perceptron and the classifier,

\begin{equation}
\label{eq:macs}
\mathrm{MACs} \;=\; N\,d \;+\; K\left(2\,d\,\lfloor d/2 \rfloor\right) \;+\; K\,d\,C,
\end{equation}

with the binned lookup contributing none. There is no term quadratic in $N$ and
no term quadratic in $K$, which is the property that makes the cost a fixed
shape set by $(N, d, K, C)$. Deployed INT8 weight storage, by contrast,
includes the full embedding table, of which only one bin per feature is read
per inference. Flash therefore exceeds the multiply-accumulate count, and this
flash-for-compute trade is quantified in Section~\ref{sec:ondevice}.

The parameter columns of Tables~\ref{tab:frontier-edge} and
\ref{tab:frontier-cic} give the trained tensor count of the deployed $R = 1$
configuration, and the flash column of Tables~\ref{tab:dev-edge} and
\ref{tab:dev-cic} gives what is actually flashed, namely the quantised weight
tensors of the modules defined in Sections~\ref{sec:method:encoder} to
\ref{sec:method:head}, including the binned embedding table, which contributes
no multiplications to \eqref{eq:macs}. The two reconcile exactly. For
Edge-IIoTset at $d = 8$ and $K = 8$, the 4,948 parameters comprise 4,912
deployed weight bytes, 27 biases, the scalar gate and the 8 per-group shift
exponents. The same accounting holds for the other deployed configurations:
15,180 against 15,120 flashed bytes at $d = 24$ on Edge-IIoTset, and 5,717
against 5,688 and 17,501 against 17,448 on CIC-IIoT-2025.

This accounting also settles where feature selection would help, a lever that
TinyML detectors such as \cite{hamdouchi2025tinyml} treat as primary. The
feature count $N$ enters \eqref{eq:macs} only through the rank-1 encoder term
$N d$, but it enters storage through both the embedding table $N \times B
\times d$ and the rank-1 direction $N \times d$. On Edge-IIoTset the
$N$-linear terms are 79.2\% of flash at $d = 8$ and 77.1\% at $d = 24$, of
which the embedding table alone is roughly seven tenths, whereas the same
terms are only 22.7\% and 14.8\% of the multiply-accumulate count, which is
dominated instead by the classifier head at 50.4\% and 32.8\%. Halving the
feature count would therefore cut the read-only footprint by close to
40\% while leaving latency largely unchanged. Feature selection is a storage
lever for this architecture rather than a compute lever, which is the opposite
of its role in models whose cost is dominated by input width, and it is not
pursued here because the binding constraint in
Section~\ref{sec:results} is latency rather than flash.

\subsection{Training protocol and device export}
\label{sec:method:training}

All neural arms share a single recipe: the Adam optimiser at a learning rate of
$3 \times 10^{-3}$, a batch size of 1024, 25 epochs, and class-weighted
cross-entropy with inverse-frequency weights, since the severe class imbalance
on all three corpora makes plain accuracy misleading. Class weighting is preferred
here to the resampling pipelines common in this literature, such as the SMOTE
and Tomek-link stage of \cite{villafranca2025edgedl}, for a reason specific to
the protocol of Section~\ref{sec:method:data}: synthesising minority examples
before a chronological split would interpolate new records from neighbours that
may lie on the other side of the partition boundary, reintroducing the
contamination the temporal split is constructed to prevent. Weighting the loss
changes the objective without manufacturing records, and so leaves the split
intact. The checkpoint retained is the
one with the best validation macro-F1, and ten seeds are reported per arm on
all three corpora. Seeds drive weight initialisation, batch order and the draw of
the group map $g$, so each reported seed uses a different random grouping and
the seed spread already absorbs topology variance.

After training, the model is quantised to INT8 by either post-training
quantisation or quantisation-aware training, and a C header is emitted
containing the quantised weight tensors, the quantile bin edges, the group map
and the per-group shift exponents. The header is compiled into firmware and
flashed to the target device, after which the device argmax is compared against
the host model argmax over 1000 replayed test flows. The criterion is exact
agreement of the predicted class index, not of the logits. All fifteen deployed
builds reported in Tables~\ref{tab:dev-edge}, \ref{tab:dev-cic} and
\ref{tab:dev-cicids} reproduced
the host prediction on every flow, at 1000/1000 (100.0\%). Reported latency is
the median of 256 inferences with the cache-cold first call excluded, since
that call runs high and would skew a plain mean. The exclusion is not
load-bearing: for the deployed Edge-IIoTset configuration the cold call costs
\SI{250.2}{\micro\second} against a warm median of \SI{243.5}{\micro\second}, so
an endpoint classifying flows sporadically rather than back to back would pay
about \SI{2.8}{\percent} more per inference. This verification
step is not a formality: it uncovered a silent export defect that was
subsequently corrected.

The deployed integer path is specified as follows. Weights are quantised to
INT8 with a single symmetric per-tensor scale for each of the embedding table,
the rank-1 direction, the two shared perceptron matrices and the classifier.
activations use an asymmetric unsigned 8-bit quantiser with an explicit zero
point at the input and again at the group state. and the biases, the residual
gate and the per-group shift exponents are retained in float. One point is
stated plainly because it bounds what the measurements establish: the deployed
kernel reproduces the arithmetic of the host quantisation-aware training loop
by dequantising each INT8 tensor to float32 and accumulating in float32, so the
device recomputes the fake-quantised forward pass the host optimised rather
than approximating it. A pure-integer kernel in the TensorFlow Lite Micro style
\cite{david2021tflm}, with int8 by int8 products accumulated into int32 and a
single fixed-point requantisation multiply, is not used here. The INT8 storage
saving reported in Tables~\ref{tab:dev-edge} and \ref{tab:dev-cic} is therefore
real and is what is flashed, and the latency reported is the true cost of the
gather-and-accumulate operator on the LX7 core, but the arithmetic itself is
float32 and the power-of-two rescale is applied as a multiplication by two to
the alpha, which is exact in binary floating point. Converting the kernel to
integer-only arithmetic is left to future work. Firmware was built with the
Arduino ESP32 core for the ESP32-S3 with the CPU clock fixed at 240\,MHz, and
latency was counted in CPU cycles on the device.

\begin{algorithm}[!ht]
\caption{ATOM training and on-device evaluation.}
\label{alg:atom}
\KwIn{train/val/test splits. embedding dim $d$, groups $K$, bins $B$, class weights $w$.}
\KwOut{test macro-F1, per-class F1, measured latency, flash bytes and parity.}
Compute per-feature quantile bin edges on the training split\;
Draw fixed random group map $g$ and freeze it\;
Initialise $E$, $V$, $W_{1}$, $W_{2}$, $W_{c}$, $b_{c}$, group scales and gate\;
\For{$\mathrm{epoch} = 1$ \KwTo $25$}{
  \ForEach{minibatch $(x, y)$}{
    Encode $e_{i}$ by \eqref{eq:encoder}\;
    Accumulate into $G_{k}$ by \eqref{eq:group}\;
    Rescale by \eqref{eq:pow2}\;
    Refine by \eqref{eq:refine}\;
    Classify by \eqref{eq:head}\;
    Step Adam on class-weighted cross-entropy\;
  }
  Keep the checkpoint with best validation macro-F1\;
}
Quantise to INT8 by PTQ or QAT and emit the C header (weights, bin edges, group map, scales)\;
Flash the ESP32-S3 and verify the device argmax against the host over 1000 flows\;
Report median latency over 256 inferences, excluding the cold call\;
\Return test macro-F1, per-class F1, measured latency, flash bytes and parity\;
\end{algorithm}

\fi

\input{results_compact}

\section{Conclusion}
\label{sec:conclusion}

ATOM-IDS replaces per-feature tokenisation with a fixed-shape grouped
accumulator for microcontroller-class intrusion detection. Binned and rank-1
feature encodings are streamed into $K$ fixed groups, rescaled by power-of-two
factors and refined by a shared residual MLP. The forward pass therefore avoids
an $N\times d$ token matrix and limits the group state to $K\times d$, independent
of feature count. This structure provides the central deployment advantage:
small activation memory without reducing every feature to a single scalar.

Under the same leakage-controlled, class-stratified chronological protocol,
the deployed $d=8$, $K=8$ model achieves macro-F1 scores of 0.8007, 0.5638 and
0.6163 on Edge-IIoTset, CIC-IIoT-2025 and CICIDS2017, respectively. Its
corresponding ESP32-S3 latencies are \SI{243.5}{\micro\second},
\SI{250.6}{\micro\second} and \SI{277.4}{\micro\second}. On Edge-IIoTset it is
$2.4\times$ faster and 7.4 macro-F1 percentage points better than the identically
deployed $h=64$ MLP, while consuming \SI{9.56}{\micro\joule} per inference at
\SI{240}{\mega\hertz}. Across the evaluated configurations, ATOM extends the
accuracy--latency frontier above the lowest-latency $h=16$ MLP point; it does
not displace that point. The tested FT-Transformer remains outside the target
arithmetic and activation-memory budgets.

These results also define the method's limits. The gather-and-accumulate kernel
is less efficient per multiply-accumulate than dense matrix multiplication on
the Xtensa LX7, and binned embeddings exchange compute for read-only storage.
Moreover, the CIC-IIoT-2025 build flags \SI{65.1}{\percent} of benign windows,
so its favourable relative cost does not establish readiness for unattended
deployment on that corpus. The reported conclusions are confined to the three
datasets, the class-stratified chronological protocol and the measured
ESP32-S3 implementation.

Future work should implement a pure-integer accumulation kernel, evaluate Arm
Cortex-M targets, calibrate operating thresholds for lower false-alarm burden
and test open-set generalisation against previously unseen attack campaigns.

%% Journal back matter: contributions, declarations, availability, references
%% and author biographies.

\section*{CRediT authorship contribution statement}

\textbf{Ashish Kalra:} Conceptualization, Methodology, Software, Validation,
Formal analysis, Investigation, Data curation, Writing -- original draft,
Visualization. \textbf{Manoj Kumar:} Conceptualization, Methodology,
Supervision, Writing -- review \& editing, Project administration.

\section*{Declaration of competing interest}

The authors declare that they have no known competing financial interests or
personal relationships that could have appeared to influence the work reported
in this paper.

\section*{Acknowledgements}

The authors thank the Department of Computer Science \& Engineering, Delhi
Technological University, for the computational resources that supported this
work.

\section*{Funding}

This research did not receive any specific grant from funding agencies in the
public, commercial or not-for-profit sectors.

\section*{Data availability}

The datasets analysed in this study are publicly available. Edge-IIoTset is
available at
\url{https://ieee-dataport.org/documents/edge-iiotset-new-comprehensive-realistic-cyber-security-dataset-iot-and-iiot-applications},
CIC IIoT 2025 (DataSense) at
\url{https://www.unb.ca/cic/datasets/iiot-dataset-2025.html}, and
CICIDS2017 at \url{https://www.unb.ca/cic/datasets/ids-2017.html}. The
preprocessing, training and device-export code supporting the findings of this
study is available from the corresponding author on reasonable request.

\bibliographystyle{IEEEtran}
\bibliography{refs}

\clearpage
\section*{Authors}

\noindent
\begin{minipage}[t]{0.25\columnwidth}
\vspace{0pt}\includegraphics[width=\linewidth]{figs/ashish_kalra.jpeg}
\end{minipage}\hfill
\begin{minipage}[t]{0.71\columnwidth}
\vspace{0pt}\textbf{Ashish Kalra} is a software engineer and a part-time PhD
research scholar in the Department of Computer Science \& Engineering at Delhi
Technological University. He received his MTech from Delhi Technological
University and his BTech from Punjab Technical University. His research
interests include security in Internet of Things (IoT) environments and
emerging machine-learning technologies. Alongside his research, he has
substantial industry experience and has held several managerial positions.
He is the corresponding author. Email:
\nolinkurl{ashishkalra_2k22phdco502@dtu.ac.in}.
\end{minipage}

\newpage
\noindent
\begin{minipage}[t]{0.25\columnwidth}
\vspace{0pt}\includegraphics[width=\linewidth]{figs/manoj_kumar.png}
\end{minipage}\hfill
\begin{minipage}[t]{0.71\columnwidth}
\vspace{0pt}\textbf{Manoj Kumar} is a professor in the Department of Computer
Science \& Engineering at Delhi Technological University. He received his
MTech from IIT Delhi and his PhD from Delhi Technological University. His
research interests include information security and natural language
processing, and his current project concerns automated article writing. He has
extensive experience in planning, budgeting and resource deployment for the
efficient accomplishment of operational goals, as well as in teaching,
training students and supervising projects. Email:
\texttt{mkumarg@dce.ac.in}.
\end{minipage}

\end{document}
